а) Произведение примитивных многочленов примитивно.

Пусть $K$ -- факториальное кольцо.
Пусть многочлены $f, g \in K[x] \setminus \{0\}$ -- примитивны. 
Докажем, что $fg$ также примитивен.

Предположим, что $fg$ не примитивен. Тогда существует неразложимый элемент
$p \in K$ такой, что $p\ |\ fg$. Но $p$ также прост в $K[x]$ (по утв 34 билета), поэтому
$p\ |\ f$ или $p\ |\ g$, но это невозможно, так как
$f$ и $g$ -- примитивны.
\bigbreak
б) Пусть $f$ -- примитивный многочлен из $\Z[x]$. 
Если $f(x)$ неприводим в $\Q[x]$, то он неприводим и в $\Z[x]$.

$\Z[x]$ факториально, поэтому если $f$ приводим над $\Z[x]$, то он
автоматически приводим и над $\Q[x]$.